% Приложения. По А.3: изходен текст на програми, формати на структури от
% данни, екранни форми, получени резултати и др.
% Приложенията НЕ влизат в препоръчителния обем на изложението (50–60 стр.).
\phantomsection
\addcontentsline{toc}{section}{ПРИЛОЖЕНИЯ}
\section*{ПРИЛОЖЕНИЯ}
\label{sec:appendix}

\subsection*{Приложение А. Пълен изходен текст на срез Button}
\addcontentsline{toc}{subsection}{Приложение А. Пълен изходен текст на срез Button}
Срезът за бутон в пълен вид --- договорът, обработчикът с двете си таблици и
платформената реализация за една среда. В Раздел~\ref{sec:implementation} са дадени само
откъсите, носещи проектните решения (листинги~\ref{lst:button_handler}
и~\ref{lst:button_ios}); тук се привежда цялото, за да е проверимо, че между откъса и
цялото няма разлика.

Показани са само \emph{относимите} откъси, а не целите файлове: пълните текстове са в
хранилището и възпроизвеждането им тук би добавило десетки страници без да добави
информация. Всеки откъс се \emph{чете} от истинския файл при компилация
(\code{\textbackslash lstinputlisting} с диапазон от редове), тоест не може да се
разминее с изходния текст --- но и не го дублира.

\lstinputlisting[caption={Договор --- \code{i\_button}: само свойства и събития, без
изобразяване}, label={lst:app_i_button}, firstline=22, lastline=32,
  basicstyle=\scriptsize\ttfamily]{../../include/maui/core/i_button.hpp}

\lstinputlisting[caption={Обявяване на обработчика --- трите таблици и създаването
на native контролата}, label={lst:app_bh_h}, firstline=199, lastline=222,
  basicstyle=\scriptsize\ttfamily]{../../include/maui/core/button_handler.hpp}

\lstinputlisting[caption={Преносима реализация --- веригата от таблици},
  label={lst:app_bh_cpp}, firstline=49, lastline=80,
  basicstyle=\scriptsize\ttfamily]{../../src/core/button_handler.cpp}

\lstinputlisting[caption={Платформена реализация (UIKit) --- създаване и едно действие},
  label={lst:app_bh_ios}, firstline=378, lastline=400,
  basicstyle=\scriptsize\ttfamily]{../../src/platform/ios/button_handler.mm}

\subsection*{Приложение Б. Публичен програмен интерфейс на системата}
\addcontentsline{toc}{subsection}{Приложение Б. Публичен програмен интерфейс на системата}
Публичните типове по слоеве, съпоставени с типовете на еталона. Съответствието е едно към
едно по покритие; имената се преобразуват по правилата на езиковия профил, а всеки заглавен
файл носи в коментар пълното име на съответстващия тип. Списъкът е генериран от първия ред на
всеки публичен заглавен файл, който носи съответствието в коментар (\code{tools/gen\_tex.py}),
и обхваща \textbf{608} типа в 19 слоя --- от \code{core} и \code{controls} до \code{layouts},
\code{graphics}, \code{hosting}, \code{bindings} и \code{xaml}.

Пълните 608 реда не се препечатват тук --- добавят над 20 страници без да добавят информация
отвъд самия списък. Табл.~\ref{tab:api-sample} показва ръчно подбрана извадка от основните
слоеве, представителна за формата и покритието на цялата таблица.

\begin{table}[H]
\centering
\caption{Представителна извадка от публичния програмен интерфейс (12 от 608 типа)}
\label{tab:api-sample}
\small
\begin{tabular}{@{}L{1.7cm}L{5.9cm}L{6.5cm}@{}}
\hline \textbf{Слой} & \textbf{Система} & \textbf{Еталон} \\ \hline
{\small core} & {\ttfamily\small maui::\allowbreak{}core::\allowbreak{}i\_\allowbreak{}element\_\allowbreak{}handler} & {\ttfamily\small Microsoft.\allowbreak{}Maui.\allowbreak{}IElement\allowbreak{}Handler} \\
{\small core} & {\ttfamily\small maui::\allowbreak{}core::\allowbreak{}i\_\allowbreak{}layout\_\allowbreak{}handler} & {\ttfamily\small Microsoft.\allowbreak{}Maui.\allowbreak{}ILayout\allowbreak{}Handler} \\
{\small core} & {\ttfamily\small maui::\allowbreak{}core::\allowbreak{}button\_\allowbreak{}handler} & {\ttfamily\small Microsoft.\allowbreak{}Maui.\allowbreak{}Handlers.\allowbreak{}Button\allowbreak{}Handler} \\
{\small controls} & {\ttfamily\small maui::\allowbreak{}controls::\allowbreak{}button} & {\ttfamily\small Microsoft.\allowbreak{}Maui.\allowbreak{}Controls.\allowbreak{}Button} \small (M2 subset) \\
{\small controls} & {\ttfamily\small maui::\allowbreak{}controls::\allowbreak{}border} & {\ttfamily\small Microsoft.\allowbreak{}Maui.\allowbreak{}Controls.\allowbreak{}Border} \\
{\small controls} & {\ttfamily\small maui::\allowbreak{}controls::\allowbreak{}absolute\_\allowbreak{}layout} & {\ttfamily\small Microsoft.\allowbreak{}Maui.\allowbreak{}Controls.\allowbreak{}Absolute\allowbreak{}Layout} \\
{\small layouts} & {\ttfamily\small maui::\allowbreak{}layouts::\allowbreak{}layout\_\allowbreak{}manager} & {\ttfamily\small Microsoft.\allowbreak{}Maui.\allowbreak{}Layouts.\allowbreak{}Layout\allowbreak{}Manager} \\
{\small layouts} & {\ttfamily\small maui::\allowbreak{}layouts::\allowbreak{}grid\_\allowbreak{}layout\_\allowbreak{}manager} & {\ttfamily\small Microsoft.\allowbreak{}Maui.\allowbreak{}Layouts.\allowbreak{}Grid\allowbreak{}Layout\allowbreak{}Manager} \\
{\small graphics} & {\ttfamily\small maui::\allowbreak{}graphics::\allowbreak{}color} & {\ttfamily\small Microsoft.\allowbreak{}Maui.\allowbreak{}Graphics.\allowbreak{}Color} \\
{\small hosting} & {\ttfamily\small maui::\allowbreak{}hosting::\allowbreak{}maui\_\allowbreak{}app\_\allowbreak{}builder} & {\ttfamily\small Microsoft.\allowbreak{}Maui.\allowbreak{}Hosting.\allowbreak{}Maui\allowbreak{}App\allowbreak{}Builder} \\
{\small bindings} & {\ttfamily\small maui::\allowbreak{}controls::\allowbreak{}binding} & {\ttfamily\small Microsoft.\allowbreak{}Maui.\allowbreak{}Controls.\allowbreak{}Binding} \\
{\small xaml} & {\ttfamily\small maui::\allowbreak{}xaml::\allowbreak{}i\_\allowbreak{}markup\_\allowbreak{}extension} & {\ttfamily\small Microsoft.\allowbreak{}Maui.\allowbreak{}Controls.\allowbreak{}Xaml.\allowbreak{}IMarkup\allowbreak{}Extension} \\ \hline
\end{tabular}
\end{table}

Пълният списък от 608 типа е генериран автоматично от изходния текст --- от първия ред на
всеки публичен заглавен файл (\code{tools/gen\_tex.py}) --- и е достъпен в хранилището на
ангажимент \code{a1d6dd9956}: заглавните файлове са на
\url{https://github.com/Alex-Tsvetanov/maui/tree/a1d6dd9956fa39126a094dcc43545b6beae5ef93/port/cpp/include/maui},
а генераторът --- на
\url{https://github.com/Alex-Tsvetanov/maui/blob/a1d6dd9956fa39126a094dcc43545b6beae5ef93/port/cpp/docs/thesis/tools/gen_tex.py}.

\subsection*{Приложение В. Съвпадение по интерфейси}
\addcontentsline{toc}{subsection}{Приложение В. Съвпадение по интерфейси}
Всеки от \boardPages{} интерфейса се оценява на четирите платформи по структурно сходство
(§\ref{subsec:scoring}) и попада в едно от три нива: съвпада, различава се незначително или
се различава съществено. \tabref{tab:heat} обобщава разпределението по платформа.

\begin{table}[H]
\centering
\caption{Разпределение на оценките по платформа}
\label{tab:heat}
\small
\begin{tabular}{@{}lccc@{}}
\hline
\textbf{Платформа} & \textbf{Съвпада (зелено)} & \textbf{Незначително (жълто)} & \textbf{Съществено (червено)} \\ \hline
iOS & \boardGreenios{} & \boardYellowios{} & \boardRedios{} \\
Mac Catalyst & \boardGreenmaccatalyst{} & \boardYellowmaccatalyst{} & \boardRedmaccatalyst{} \\
Android & \boardGreenandroid{} & \boardYellowandroid{} & \boardRedandroid{} \\
Windows & \boardGreenwindows{} & \boardYellowwindows{} & \boardRedwindows{} \\ \hline
\end{tabular}
\end{table}

Таблицата показва входния път \emph{по код}: по него няма нито един интерфейс със съществено
разминаване на нито една от четирите платформи. По входния път \emph{от описанието} остава едно
изключение --- iOS, \code{selection\_synchronization} --- разгледано в §\ref{subsec:visual}
(\tabref{tab:worst}) заедно с констатацията защо не е принудено до зелено.

Разбивката по \textbf{отделен интерфейс} --- каквато носеше предишният вид на тази таблица,
ред по ред --- не се препечатва тук: строи се възпроизводимо от \code{comparison.json}
(\code{tools/gen\_tex.py}). Изходните данни за всеки от \boardPages{} интерфейса, по двата
входни пътя (код / описание), са версионирани на същия ангажимент:
\url{https://github.com/Alex-Tsvetanov/maui/blob/a1d6dd9956fa39126a094dcc43545b6beae5ef93/port/cpp/docs/comparison/comparison.json};
четим изглед със статус по интерфейс и платформа --- в
\url{https://github.com/Alex-Tsvetanov/maui/blob/a1d6dd9956fa39126a094dcc43545b6beae5ef93/port/cpp/docs/comparison/README.md}.

\subsubsection*{Най-слабо съвпадащите интерфейси --- съпоставки}

Всеки интерфейс от \tabref{tab:worst} е показан тук един до друг с еталона, в темата, която
дърпа оценката надолу --- показването на светлия кадър за разминаване, което съществува само
в тъмната тема, не илюстрира нищо. Измерените стойности стоят в заглавието на дясната колона
и в легендата, така че твърдението в текста може да бъде сверено с изображението, вместо да
се приема на доверие.

\input{generated/worst_figures}

\subsection*{Приложение Г. Екранни форми}
\addcontentsline{toc}{subsection}{Приложение Г. Екранни форми}
Съпоставки за интерфейси, които не са включени в Раздел~\ref{sec:results}. Всички са
символни връзки към заснеманията от борда, тоест същите файлове, които са били оценени.

\evidencepair{maccatalyst}{alignment}{light}{Подравняване --- проверява разположението без
собствено изобразяване на контролите.}

\evidencepair{maccatalyst}{clip}{dark}{Отсичане, тъмна тема --- класът, при който
половинпикселното свиване на контура се проявява (§\ref{subsec:taxonomy}).}

\evidencepair{ios}{border_playground}{light}{Рамки и радиуси --- интерфейсът, върху който
е измерена промяната в размера на превключвателя (§\ref{subsec:port_strategy}).}

\evidencepair{ios}{grid}{light}{Решетка --- разпределяне на пропорционални измерения.}

\evidencepair{ios}{entry}{dark}{Поле за въвеждане, тъмна тема --- интерфейсът, върху който
се прояви размерът на превключвателя (§\ref{subsec:port_strategy}).}

\evidencepair{android}{alerts}{light}{Съобщения --- native диалози на трета среда.}

\noindent\small За анимираните интерфейси (\gifPageCount{} на брой) съпоставката би
показала първи кадър; тук такива не са включени.\normalsize
